\section{Implementation and On-Road Deployment}
\label{sec:road}

\noindent\textbf{Hardware.} The rig is built from commodity parts (Fig.~\ref{fig:teaser}): a mid-range Android phone (OnePlus Nord N10, US\$150), an NVIDIA Jetson Orin Nano developer kit (US\$250), and US\$100 of mount, 12\,V inverter and cable, US\$500 in total. Nothing connects to the vehicle's controls or its bus, and no server takes part in computing an advisory: the only network traffic is the feed query.

\noindent\textbf{Software.} The phone captures the camera, GPS and inertial unit, queries the feed once a minute, and forwards everything raw over USB; the Jetson runs detection (YOLOv8n at 640\,px), tracking, distance and density estimation, both advisory layers, the safety filter, and a scheduler that sets each sensor's sampling rate. Interpretation runs on the Jetson because computing on the phone as well would worsen the thermal problem below.

\noindent\textbf{Drives.} We drove the rig on 2026-09-08 on public roads in central New Jersey. Eight runs recorded 22{,}929 ticks over 152.8 minutes, with a valid GPS fix on 99.1\% of them, and 123 feed responses, every one returning HTTP 200. In two runs the phone applied the sampling rates the scheduler sent, and in six it only logged them. No driver acted on an advisory: that needs institutional review board approval, which this project does not hold, and one instrumented vehicle cannot change how traffic flows, so the traffic results are simulated (\S~\ref{sec:eval}).

\noindent\textbf{Computation is not the bottleneck; heat is.} On the Jetson the path from detection to the decoded advisory takes a median of 19.8\,ms, and the phone-to-Jetson link has a 95th-percentile latency of 116.19\,ms over USB against 215.63\,ms over a network link, which is why the rig uses USB. The failures were physical. On the first drive the phone reached Android's severe thermal tier on a windscreen mount in sunlight and the platform closed the application; a failed cable and a hotspot shut off by sunlight cost two further drives.

\noindent\textbf{From recorded frames to the simulated camera.} We ran the rig's own detector, tracker and distance code on the 18,818 frames that six drives indexed against the other sensors; the simulated camera uses the error tables of the four drives recorded at five frames per second. Measured against a second pass at 1,280\,px, which sees farther, the detection rate falls from 0.95 within 10\,m to 0.53 at 40 to 50\,m and 0.32 at 50 to 60\,m, which sets the range at 60\,m. We compute distance from each box's width and an assumed vehicle width. Distance carries the frame-to-frame noise we measured, from 0.05\,m at 10\,m to 1.1\,m at 50 to 60\,m, and a per-vehicle offset with a spread of 31\%, because a width-based distance is wrong by as much as that vehicle's width differs from the assumed one. The simulated camera sees its own and the adjacent lanes, with nearer vehicles hiding those behind them. The phone's focal length was not recorded, so the absolute distance scale is assumed (\S~\ref{sec:sens} varies it).
